\renewcommand{\reportmonth}{Januari}
\renewcommand{\signaturedate}{1 Februari 2026}
\newpage
\TabPositions{4.5cm}

\begin{center}
    \textbf{LOG BOOK \\ MAGANG INDUSTRI}
\end{center}

\raggedright%
Program Studi           \tab: {\studyprogram} \\
Semester/Th. Akademik   \tab: {\semester} / {\academicperiod} \\
NIM/Nama Mahasiswa      \tab: {\authornim} / {\authorname} \\
Perusahaan              \tab: {\company} \\
Job Title               \tab: {\jobposition} \\
Departemen              \tab: {\jobdepartment} \\
Mentor                  \tab: {\mentor} \\
Bulan                   \tab: {\reportmonth} {\reportyear} \\
\vspace{12pt}

\begin{logtable}

    % =============================================================
    % Minggu Ke-1  (12 Januari - 18 Januari 2026)
    % =============================================================
    \logentry{1.}
    {
        Senin \linebreak 12 Jan 2026 \linebreak
        08:00--17:00 \linebreak
        WFO
    }{
        \begin{loglist}
            \item Memahami aturan, budaya kerja, serta kebijakan perusahaan melalui kegiatan
            Orientasi Karyawan Baru (OKB) di \textit{Head Office}.
            \item Mempelajari arsitektur dan \textit{framework} yang digunakan dalam
            proyek-proyek perusahaan.
        \end{loglist}
    }{
        \begin{logwork}
            \item Mengikuti OKB di \textit{Ruang Meeting} bersama karyawan dan peserta magang
            baru di \textit{Head Office}.
            \item Melakukan \textit{cloning repository} GitHub
            \texttt{healthcare-intelligence-app}, yaitu \textit{repository} khusus
            \textit{Front-End} untuk proyek HealthyOne dan Canvasser.
        \end{logwork}
    }

    \logentry{2.}
    {
        Selasa \linebreak 13 Jan 2026 \linebreak
        08:00--17:00 \linebreak
        WFO
    }{
        \begin{loglist}
            \item Mempelajari \textit{workflow} dan cara penggunaan \textit{framework} Flutter
            yang diterapkan perusahaan.
            \item Memulai pengerjaan Task \#11858 (\textit{IMPROVEMENT}): Migrasi beberapa
            \textit{widget} ke Flutter versi terbaru.
        \end{loglist}
    }{
        \begin{logwork}
            \item Berlatih membuat \textit{widget} dasar seperti \texttt{table, button,} dan
            \texttt{tabs} menggunakan \textit{local development environment}.
            \item Melakukan migrasi kode \textit{widget} dari Flutter versi sebelumnya ke Flutter
            versi terbaru, mempelajari arsitektur MVC, serta menyelesaikan
            \textit{dependency issues} yang ditemukan.
        \end{logwork}
    }

    \logentry{3.}
    {
        Rabu \linebreak 14 Jan 2026 \linebreak
        08:00--17:00 \linebreak
        WFO
    }{
        \begin{loglist}
            \item Mendapatkan Task \#11858 (\textit{IMPROVEMENT}): Migrasi beberapa
            \textit{widget} ke Flutter versi terbaru.
        \end{loglist}
    }{
        \begin{logwork}
            \item Melakukan \textit{update} dan \textit{refactor code} untuk kode-kode
            \textit{widget} yang mengalami \textit{deprecated packages}.
        \end{logwork}
    }

    \logentry{4.}
    {
        Kamis \linebreak 15 Jan 2026 \linebreak
        08:12--17:12 \linebreak
        WFO
    }{
        \begin{loglist}
            \item Melanjutkan Task \#11858
        \end{loglist}
    }{
        \begin{logwork}
            \item Melakukan perbaikan pada \textit{widgets} yang memiliki \textit{bug}.
        \end{logwork}
    }

    \logevaluate{Dokumentasi Hasil Pekerjaan: Layar Latihan untuk Membuat \textit{Widget}}{
        \noindent\includegraphics[width=\linewidth]{3.4-latihan.png}
    }

    % =============================================================
    % Minggu Ke-2  (19 Januari - 25 Januari 2026)
    % =============================================================

    \logentry{5.}
    {
        Senin \linebreak 19 Jan 2026 \linebreak
        07:58--17:10 \linebreak
        WFO
    }{
        \begin{loglist}
            \item Melanjutkan Task \#11858
        \end{loglist}
    }{
        \begin{logwork}
            \item Melakukan migrasi \textit{widgets} ke Flutter versi terbaru dan memperbaiki
            \textit{bug}.
        \end{logwork}
    }

    \logentry{6.}
    {
        Selasa \linebreak 20 Jan 2026 \linebreak
        08:12--17:15 \linebreak
        WFO
    }{
        \begin{loglist}
            \item Melanjutkan Task \#11858
        \end{loglist}
    }{
        \begin{logwork}
            \item Melakukan \textit{refactor} pada beberapa \textit{widgets} agar sesuai dengan
            arsitektur MVC.
        \end{logwork}
    }

    \logentry{7.}
    {
        Rabu \linebreak 21 Jan 2026 \linebreak
        08:19--17:21 \linebreak
        WFO
    }{
        \begin{loglist}
            \item Melanjutkan Task \#11858
        \end{loglist}
    }{
        \begin{logwork}
            \item Melakukan \textit{refactor} pada beberapa \textit{widgets}
            \texttt{(code\_snippet, image\_field)} agar sesuai dengan arsitektur MVC.
        \end{logwork}
    }

    \logentry{8.}
    {
        Kamis \linebreak 22 Jan 2026 \linebreak
        08:12--17:33 \linebreak
        WFO
    }{
        \begin{loglist}
            \item Melanjutkan Task \#11858
        \end{loglist}
    }{
        \begin{logwork}
            \item Melakukan \textit{refactor} pada beberapa \textit{widgets} \texttt{(map)} agar
            sesuai dengan arsitektur MVC.
        \end{logwork}
    }

    \logentry{9.}
    {
        Jumat \linebreak 23 Jan 2026 \linebreak
        07:58--17:13 \linebreak
        WFO
    }{
        \begin{loglist}
            \item Melanjutkan Task \#11858
        \end{loglist}
    }{
        \begin{logwork}
            \item Melakukan \textit{refactor} pada beberapa \textit{widgets}
            \texttt{(text\_field)} agar sesuai dengan arsitektur MVC.
        \end{logwork}
    }

    \logevaluate{Dokumentasi Hasil Pekerjaan: Migrasi kode \textit{widgets}}{
        \noindent\includegraphics[width=\linewidth]{3.4-migrasi-3.png}
    }

    % =============================================================
    % Minggu Ke-3  (26 Januari - 1 Februari 2026)
    % =============================================================
    \logentry{10.}
    {
        Senin \linebreak 26 Jan 2026 \linebreak
        07:46--17:17 \linebreak
        WFO
    }{
        \begin{loglist}
            \item Mendapat Task \#10319 (\textit{IMPROVEMENT}): Buat agar \texttt{flexible
                dropdown} bisa \textit{multiple choice}.
        \end{loglist}
    }{
        \begin{logwork}
            \item (\#10319) Memodifikasi \textit{text handling} agar \textit{input flexible dropdown} bisa menerima beberapa pilihan sekaligus dan bisa \textit{update state} dengan baik.
        \end{logwork}
    }

    \logentry{11.}
    {
        Selasa \linebreak 27 Jan 2026 \linebreak
        08:00--17:14 \linebreak
        WFO
    }{
        \begin{loglist}
            \item Mempelajari \textit{framework} React.
            \item Mempelajari arahan terkait alur kerja perusahaan (penggunaan
            \textit{OpenProject} dan sistem KPI).
        \end{loglist}
    }{
        \begin{logwork}
            \item Mempelajari React secara otodidak dari tutorial YouTube dan
            \textit{roadmap.sh}.
            \item Arahan terkait alur kerja diberikan secara personal oleh Stephanie Lirungan,
            \textit{Deputy Project Manager} (DPM).
        \end{logwork}
    }

    \logentry{12.}
    {
        Rabu \linebreak 28 Jan 2026 \linebreak
        07:58--17:38 \linebreak
        WFO
    }{
        \begin{loglist}
            \item Mendapat Task \#10120 (\textit{IMPROVEMENT}): Buat agar \textit{widget}
            \texttt{template form} bisa memuat \textit{widget} \texttt{accordion} dengan
            \textit{subtype column/row}.
        \end{loglist}
    }{
        \begin{logwork}
            \item (\#10120) Menambahkan \textit{parser} untuk \texttt{accordion} di dalam \texttt{template.form\_dart} dan modelnya.
        \end{logwork}
    }

    \logentry{13.}
    {
        Kamis \linebreak 29 Jan 2026 \linebreak
        07:56--17:26 \linebreak
        WFO
    }{
        \begin{loglist}
            \item Mendapat Task \#10074 (\textit{IMPROVEMENT}): Buat agar \textit{widget}
            \texttt{multiple choice (checkbox)} bisa diklik di area teksnya.
        \end{loglist}
    }{
        \begin{logwork}
            \item (\#10074) Menambahkan \textit{InkWell} dalam area teks \texttt{checkbox} agar bisa diklik.
        \end{logwork}
    }

    \logentry{14.}
    {
        Jumat \linebreak 30 Jan 2026 \linebreak
        07:58--17:53 \linebreak
        WFO
    }{
        \begin{loglist}
            \item Mendapat Task \#10735 (\textit{IMPROVEMENT}): Buat agar \textit{styling} teks
            di \textit{popup AI Scribe} 1.5x lebih besar.
        \end{loglist}
    }{
        \begin{logwork}
            \item (\#10735) Menambahkan \textit{textStyle: size} 1.5x dalam teks hasil \textit{AI Scribe}.
        \end{logwork}
    }

    \logevaluate{Dokumentasi Hasil Pekerjaan: Migrasi kode \textit{widgets}}{
        \noindent\includegraphics[width=\linewidth]{3.4-migrasi-3.png}
    }

    \logevaluate{Analisis dan Pembahasan}{
        Kegiatan magang industri pada bulan Januari 2026 diawali dengan
        Orientasi Karyawan Baru (OKB) di \textit{Head Office}, pengenalan arsitektur
        proyek, serta persiapan \textit{environment} melalui \textit{cloning
            repository healthcare-intelligence-app}. Setelah \textit{onboarding}, pengerjaan
        difokuskan pada Task \#11858 yang mencakup migrasi seluruh \textit{widget} ke
        Flutter versi terbaru sekaligus \textit{refactor} agar sesuai dengan arsitektur
        MVC.\ Tugas ini melibatkan penanganan
        \textit{deprecated packages}, perbaikan \textit{bug} akibat perubahan API, serta
        restrukturisasi \textit{widget} seperti \textit{code\_snippet, image\_field, map,} dan
        \textit{text\_field}.

        Di akhir bulan, pengerjaan beralih ke beberapa task \textit{improvement} dengan
        cakupan yang lebih spesifik, yaitu Task \#10319 untuk mendukung \textit{input
            multiple choice} pada \textit{flexible dropdown}, Task \#10120 untuk
        menambahkan dukungan \textit{accordion} pada \textit{template form}, Task
        \#10074 untuk memperluas area klik pada \textit{widget checkbox}, serta Task
        \#10735 untuk menyesuaikan ukuran teks pada \textit{popup} \textit{AI Scribe}.
        Selain itu, dilakukan pula pengenalan \textit{framework} React melalui
        pembelajaran mandiri serta pemahaman alur kerja perusahaan melalui arahan dari
        \textit{Deputy Project Manager}. }

    \logevaluate{Capaian yang Telah Diselesaikan}{
        \begin{itemize}[noitemsep, topsep=2pt, leftmargin=12pt, label=\textbullet]
            \item Menyelesaikan Task \#11858: migrasi seluruh \textit{widget} ke Flutter versi
                  terbaru, \textit{refactor widget code\_snippet, image\_field, map,} dan
                  \textit{text\_field} agar sesuai dengan arsitektur MVC.
            \item Menyelesaikan Task \#10319: modifikasi \textit{flexible dropdown} untuk
                  mendukung \textit{multiple choice} dengan \textit{state management} yang
                  berfungsi dengan benar.
            \item Menyelesaikan Task \#10120: penambahan \textit{parser accordion} beserta model
                  pada \textit{widget template form}.
            \item Menyelesaikan Task \#10074: perluasan area klik \textit{checkbox} pada
                  \textit{widget multiple choice} menggunakan komponen \textit{InkWell}.
            \item Menyelesaikan Task \#10735: penerapan konfigurasi \textit{textStyle} 1.5x pada
                  teks hasil \textit{output} di \textit{popup AI Scribe}.
        \end{itemize}
    }

    \logevaluate{Kendala}{
        Proses migrasi dan \textit{refactor widget} ke Flutter versi terbaru memerlukan waktu yang cukup panjang karena banyaknya \textit{deprecated packages} serta perbedaan API yang signifikan antara versi lama dan baru, sehingga setiap \textit{widget} perlu ditinjau dan disesuaikan secara individual.
    }

    \logevaluate{Solusi}{
        Mempelajari dokumentasi resmi Flutter dan \textit{changelog} versi terbaru secara bertahap untuk memahami perubahan API, serta mengerjakan \textit{refactor} secara sistematis \textit{widget} per \textit{widget} agar setiap perubahan dapat diverifikasi fungsionalitasnya sebelum lanjut ke \textit{widget} berikutnya.
    }

\end{logtable}
% \footnotesize
\noindent
% Definisi Environment Tabel
\newenvironment{tasktable}
{
    \begin{longtable}{
        | @{\hspace{0pt}} >{\centering\arraybackslash}p{1cm} @{\hspace{0pt}} 
        | @{\hspace{0pt}} >{\centering\arraybackslash}p{3.1cm} @{\hspace{0pt}} 
        | @{\hspace{0pt}} >{\centering\arraybackslash}p{3.1cm} @{\hspace{0pt}} 
        | @{\hspace{0pt}} >{\centering\arraybackslash}p{3.1cm} @{\hspace{0pt}} 
        | >{\raggedright\arraybackslash}p{3.1cm} |}
        \hline
        \textbf{No.} & \textbf{Jenis} & \textbf{Tanggal Pengerjaan} & \textbf{Requester} & \textbf{Detail} \\ \hline
        \endfirsthead
        
        \hline
        \textbf{No.} & \textbf{Jenis} & \textbf{Tanggal Pengerjaan} & \textbf{Requester} & \textbf{Detail} \\ \hline
        \endhead
}
{ 
    \end{longtable} 
}

% Command untuk input data sesuai gambar
% #1: No
% #2: Judul/ID Tugas (Baris panjang)
% #3: Jenis (IMPROVE, dll)
% #4: Tanggal
% #5: Requester
% #6: Detail (Lorem Ipsum)
\newcommand{\taskentry}[6]{
    \multirow{2}{*}{#1} & \multicolumn{4}{l|}{#2} \\ \cline{2-5}
    & #3 & #4 & #5 & #6 \\ \hline
}
\begin{flushright}
    {\city}, {\signaturedate}\\
    Mengetahui,\\
    Mentor {\jenisMBKM}\\[2cm]
    {\mentor}\\
\end{flushright}